\begin{proposition}[A finite positivity projection on three peers]
\label{prop:c2threepositivity}
For $n=3$ and $c=\ell=2$, let $\mathcal P$ be the family in
which every cache has at most one positive-age descriptor.
Every peer's cache lies in this family after its first
participation in an exchange, and this property is permanent.
Consequently, from every weakly connected initial state with
arbitrary finite ages,
\[
 \Pr_x\{X_m\notin\mathcal P\}\le3(2/3)^m.
\]
On $\mathcal P$, replacing every positive numerical age by one
is an exact Markov projection, with $487$ weakly connected
projected states.
\end{proposition}
\begin{proof}
Write the three distinct peers as $p,q,r$, and consider
$p\to q$. After mandatory removal of $q$, the initiator can
retain only identity $r$: the reply contains only $p$ and $r$,
and $p$ is self-filtered. Thus its resulting cache has at most
one entry, regardless of the initial ages. The responder
retains protected fresh $p^0$ and can have at most one other
entry. Hence both updated caches have at most one positive
entry. Untouched caches are unchanged, proving permanence.
If each peer has initiated at least once, all three caches
satisfy the property. The independent uniform initiators and
a union bound give the displayed probability estimate.

Within this family a positive descriptor is the uniquely
oldest entry whenever a second entry is present; a singleton
descriptor is also removed on initiation. Thus no previously
positive active descriptor is ever forwarded as an extra.
An extra had age zero and is sent at age one. Fresh entries
have age zero; copies preserve positivity; and minimum merges
have age zero exactly when at least one input has age zero.
Partner ties occur exactly between two zero entries.
Therefore transition probabilities of identities and
zero/positive flags depend only on those flags, proving the
projection assertion. Each peer has eight possible projected
rows: the empty row, four singleton rows, and the three
two-entry rows with at most one positive entry. Filtering
the $8^3$ triples by weak connectivity yields $487$ states.
\end{proof}

\begin{proposition}[A certified ordinary sampling bound on three peers]
\label{prop:c2threeordinarysample}
For $n=3,c=\ell=2$, the conditional single-identifier sampling
error is at most $\varepsilon$, for every weakly connected
initial state with arbitrary finite ages and every observation
time at least
\[
 \left\lceil\frac{\log(9/\varepsilon)}{\log(3/2)}\right\rceil
 +8\left\lceil\log_2(3/\varepsilon)\right\rceil,
 \qquad 0<\varepsilon<1.
\]
This bound concerns ordinary independent uniform activations,
not a prescribed controlled execution. Its finite operator
certificate applies specifically to three peers.
\end{proposition}
\begin{proof}
Relabel the observer as $0$, and let $\sigma$ swap the other
two labels. Let $P$ be the $487$-state positivity-projection
kernel of Proposition~\ref{prop:c2threepositivity}.
The standard-library checker
\texttt{verify\_c2\_three\_peer\_antisymmetric.py}
enumerates all its transitions exactly. Writing $A=6P$, an
integer matrix, it verifies equivariance and the exact
finite inequalities
\[
 \max_z\frac12\sum_w
   \left|P^8(z,w)-P^8(\sigma z,w)\right|
   =\frac{34447}{69984}<\frac12,
 \qquad
 \min_z(P^4\mathbf1_{\{|C_0|>0\}})(z)=\frac59.
\]
These are exhaustive integer calculations, without age
truncation or floating-point inequalities: the maximum
row-difference numerator is $1653456$ over denominator
$2\cdot6^8=3359232$.

For every $\sigma$-antisymmetric function $f$,
equivariance gives
\[
 |P^8f(z)|
 =\left|\frac12\sum_w
       (P^8(z,w)-P^8(\sigma z,w))f(w)\right|
 \le\frac12\|f\|_\infty.
\]
The kernel preserves antisymmetry, so iteration and
supremum-norm contraction imply
$\|P^sf\|_\infty\le2^{-\lfloor s/8\rfloor}\|f\|_\infty$.
Choose $f(z)$ as the conditional-on-state output probability
of identity $1$ minus that of identity $2$, including the
value zero at an empty cache. This function is antisymmetric
and bounded by one. Its expected value is the difference of
the two unconditional query probabilities.
The availability certificate also persists at all $s\ge4$,
since a Markov kernel cannot decrease a function's minimum.

After $m$ initial activations, write $b=3(2/3)^m<1$.
On the event of entering $\mathcal P$, the remaining projected
process obeys these bounds; its complement has probability
at most $b$. Thus at time $m+s$, for $s\ge4$, the difference
between unconditional query probabilities is at most
$b+2^{-\lfloor s/8\rfloor}$, and availability is at least
$(1-b)5/9$. For two target identities the conditional
total-variation error is their unconditional probability
difference divided by twice availability. It is therefore
at most
\[
 \frac9{10}\,
 \frac{b+2^{-\lfloor s/8\rfloor}}{1-b}.
\]
Taking $m=\lceil\log(9/\varepsilon)/\log(3/2)\rceil$ and
$s=8\lceil\log_2(3/\varepsilon)\rceil$ makes both numerator
terms at most $\varepsilon/3$. Since $\varepsilon<1$, the
displayed error is at most $0.9\varepsilon<\varepsilon$.
The same estimates apply at every later observation time.
\end{proof}
